\section{Canonical faithful completion and the hypergraph alternative}
\label{sec:canonical-completion}

The preceding calculus decides whether one proposed pair projection is
faithful.  We now solve the inverse problem: given the literal full
incidence, when does any pair-only rectangular representation exist, and
what is its unique least relation?  The answer depends on whether one asks
to preserve the parent form or only the terminal atomic field.

\subsection{Two canonical active incidences}

Continue with the finite data of \cref{sec:incidence-calculus}.  Because the
structural support is fixed before \(z\), the following sets are also
coefficient-independent.

\begin{definition}[Pre-terminal and terminal active incidences]
\label{def:completion-active-incidences}
For a full incidence \(\cU\subseteq\Omega\), define
\begin{align}
 \cU^{\mathrm{pre},+}
 &:=\{u\in\cU:c_u\ne0\},
 \label{eq:completion-pre-active}\\
 \cU_h^+
 &:=\{u\in\cU:t_h(u)c_u\ne0\}.
 \label{eq:completion-terminal-active}
\end{align}
\end{definition}

The second set is contained in the first.  It may depend on the fixed
physical shift and terminal support, but not on the row coefficient vector.
It must not be substituted for \(\cU^{\mathrm{pre},+}\) in an upstream
parent form: that substitution would make terminal activation true by
definition.

\begin{proposition}[Unique minimal coordinate supports]
\label{prop:completion-minimal-coordinate-supports}
Let \(V\subseteq\cU\).
\begin{enumerate}[label=\textup{(\roman*)},leftmargin=3em]
 \item The following are equivalent:
       \begin{equation}
        C_V=C_{\cU},
        \qquad
        P_V=P_{\cU}\text{ on }Z,
        \qquad
        \cU^{\mathrm{pre},+}\subseteq V.
        \label{eq:completion-pre-minimal-equivalence}
       \end{equation}
       Consequently \(\cU^{\mathrm{pre},+}\) is the unique least
       coordinate support preserving the pre-terminal synthesis and form.
 \item The following are equivalent:
       \begin{equation}
        D_hC_V=D_hC_{\cU},
        \qquad
        A_{h,V}=A_{h,\cU}\text{ on }Z,
        \qquad
        \cU_h^+\subseteq V.
        \label{eq:completion-terminal-minimal-equivalence}
       \end{equation}
       Consequently \(\cU_h^+\) is the unique least coordinate support
       preserving the terminal atomic synthesis and form.
\end{enumerate}
\end{proposition}

\begin{proof}
If \(V\) contains the stated active set, every omitted coordinate has zero
coefficient in the relevant synthesis.  Conversely, if
\(u_0\in\cU^{\mathrm{pre},+}\setminus V\), then, at
\(z=e_{m(u_0)}\),
\[
 P_{\cU}(z)-P_V(z)\ge |c_{u_0}|^2>0.
\]
The terminal statement follows identically with
\(|t_h(u_0)c_{u_0}|^2\).  Thus every preserving support contains the
corresponding active set.
\end{proof}

There need not be a unique least coordinate support preserving only the
grouped operator.  Exact cancellations inside a grouped fiber can make
different subsets synthesize the same output.  The terminal atomic support
is canonical because it is defined before that cancellation.

\subsection{Active rectangular closure}

For \(I\subseteq\cM\times\cR\), define
\begin{align}
 \operatorname{PRect}(I)
 &:=\{u\in\Rect(I):c_u\ne0\},
 \label{eq:completion-PRect}\\
 \TRect_h(I)
 &:=\{u\in\Rect(I):t_h(u)c_u\ne0\}.
 \label{eq:completion-TRect}
\end{align}
Both operations are monotone in \(I\).

\begin{definition}[Atomic-faithful pair relation]
\label{def:completion-atomic-faithful}
A pair relation \(I\) is
\begin{enumerate}[label=\textup{(\alph*)},leftmargin=2.7em]
 \item \emph{pre-terminal faithful} for \(\cU\) if
       \begin{equation}
        \operatorname{PRect}(I)=\cU^{\mathrm{pre},+};
        \label{eq:completion-pre-faithful}
       \end{equation}
 \item \emph{terminal-atomic faithful} for \((\cU,h)\) if
       \begin{equation}
        \TRect_h(I)=\cU_h^+.
        \label{eq:completion-terminal-faithful}
       \end{equation}
\end{enumerate}
\end{definition}

These definitions demand both completeness and absence of active spurious
crosses.  They are stronger than equality after grouping and are the
appropriate notions for raw exposure forms.

\begin{theorem}[Unique least terminal-faithful pair relation]
\label{thm:completion-least-terminal-faithful-relation}
Put
\begin{equation}
 I_h^{\min}:=\pr_{m,r}\cU_h^+.
 \label{eq:completion-Ih-min}
\end{equation}
There exists a terminal-atomic faithful pair relation for \((\cU,h)\) if
and only if
\begin{equation}
 \boxed{\TRect_h(I_h^{\min})=\cU_h^+.}
 \label{eq:completion-terminal-rectangularity-criterion}
\end{equation}
When this condition holds, \(I_h^{\min}\) is the unique least faithful
relation.  More precisely, \(I\) is terminal-atomic faithful if and only if
\begin{equation}
 I_h^{\min}\subseteq I
 \quad\text{and}\quad
 \TRect_h(I\setminus I_h^{\min})=\varnothing.
 \label{eq:completion-all-terminal-faithful-relations}
\end{equation}
\end{theorem}

\begin{proof}
Every \(u\in\cU_h^+\) has its pair in \(I_h^{\min}\), hence
\begin{equation}
 \cU_h^+\subseteq\TRect_h(I_h^{\min}).
 \label{eq:completion-active-contained-in-min-closure}
\end{equation}
If a faithful relation \(I\) exists, it contains the pair of every
\(u\in\cU_h^+\), so \(I_h^{\min}\subseteq I\).  Monotonicity gives
\[
 \cU_h^+
 \subseteq\TRect_h(I_h^{\min})
 \subseteq\TRect_h(I)
 =\cU_h^+,
\]
which proves \eqref{eq:completion-terminal-rectangularity-criterion}.
Conversely, that condition says exactly that \(I_h^{\min}\) is faithful,
and every faithful relation contains it.

Assume the criterion holds.  If \(I\) is faithful, any active atom created
by a pair in \(I\setminus I_h^{\min}\) has a pair absent from the
projection of \(\cU_h^+\), and so lies outside \(\cU_h^+\), a
contradiction.  This proves necessity of the second condition in
\eqref{eq:completion-all-terminal-faithful-relations}.  The two conditions
are sufficient because the first retains all of
\(\TRect_h(I_h^{\min})=\cU_h^+\), while the second adds no active atom.
\end{proof}

\begin{corollary}[Unique least pre-terminal faithful relation]
\label{cor:completion-least-preterminal-faithful-relation}
Put
\begin{equation}
 I_{\mathrm{pre}}^{\min}
 :=\pr_{m,r}\cU^{\mathrm{pre},+}.
 \label{eq:completion-Ipre-min}
\end{equation}
A pre-terminal faithful pair relation exists if and only if
\begin{equation}
 \boxed{
 \operatorname{PRect}(I_{\mathrm{pre}}^{\min})
 =\cU^{\mathrm{pre},+}.}
 \label{eq:completion-pre-rectangularity-criterion}
\end{equation}
When it exists, \(I_{\mathrm{pre}}^{\min}\) is the unique least one.
\end{corollary}

\begin{proof}
Repeat the proof of
\cref{thm:completion-least-terminal-faithful-relation}, replacing
\(\TRect_h\) by \(\operatorname{PRect}\) and \(\cU_h^+\) by
\(\cU^{\mathrm{pre},+}\).
\end{proof}

Condition \eqref{eq:completion-pre-rectangularity-criterion} is stronger
than \eqref{eq:completion-terminal-rectangularity-criterion}.  A crossed
coordinate can carry nonzero parent energy and still be terminal-null.
Thus exact terminal factorization through a pair relation does not justify
using the rectangular lift's parent form as the upstream parent form.

\subsection{Pair projection is not full provenance}

\begin{example}[The same pair projection supports different incidences]
\label{ex:completion-same-pair-different-incidence}
Take one row \(m\), two cofactors \(r_1,r_2\), and two base labels
\(b_1,b_2\in\cJ\times\Gamma\).  In the resulting four-point rectangle let
\begin{align}
 \cU_{\mathrm{diag}}
 &:=\{(m,r_1,b_1),(m,r_2,b_2)\},
 \label{eq:completion-diagonal-incidence}\\
 \cU_{\mathrm{cross}}
 &:=\{(m,r_1,b_2),(m,r_2,b_1)\},
 \label{eq:completion-cross-incidence}\\
 \cU_{\mathrm{full}}
 &:=\{m\}\times\{r_1,r_2\}\times\{b_1,b_2\}.
 \label{eq:completion-full-incidence}
\end{align}
All three pair projections equal
\(\{(m,r_1),(m,r_2)\}\).  Let every point of the ambient rectangle be
terminal-active and let the output key resolve individual atoms.  The
projected relation is faithful for \(\cU_{\mathrm{full}}\), but no pair
relation is terminal-atomic faithful for either
\(\cU_{\mathrm{diag}}\) or \(\cU_{\mathrm{cross}}\): retaining both actual
pairs necessarily creates the other two active crosses.

The atom-resolved operators are different even if all coefficients have
modulus one.  To make the scalar raw forms differ as well, assign terminal
coefficient modulus one to the diagonal points and modulus two to the
cross points.  At \(z=e_m\), the diagonal and cross forms are then \(2\)
and \(8\), while their pair projections remain equal.
\end{example}

\begin{example}[Grouped equality can conceal a nonfaithful rectangle]
\label{ex:completion-grouped-concealment}
Suppose
\(\Rect(I_h^{\min})\setminus\cU_h^+\) contains two active crosses with the
same row and grouped output key, having terminal coefficients \(1\) and
\(-1\).  The rectangular lift is not terminal-atomic faithful, yet the
spurious columns cancel in the grouped operator.  Hence grouped equality
does not certify an exposure identity.  Under
\(q=(\gamma,d_mr,j)\), this cancellation is excluded by
\cref{cor:incidence-physical-key-singleton}.
\end{example}

These are finite-incidence countermodels, not asserted arithmetic
counterexamples inside the actual TPC carrier.  Their exact conclusion is
that hypotheses recording only the pair projection cannot reconstruct the
full incidence.

\subsection{When the full incidence hypergraph is mandatory}

\begin{theorem}[Sharp hypergraph-necessity criterion]
\label{thm:completion-hypergraph-necessity}
If
\begin{equation}
 \TRect_h(I_h^{\min})\setminus\cU_h^+\ne\varnothing,
 \label{eq:completion-active-spurious-obstruction}
\end{equation}
then no pair relation is simultaneously complete for every atom of
\(\cU_h^+\) and free of terminal-active spurious coordinates.  The full
incidence hypergraph
\[
 \cU_h^+\subseteq\cM\times\cR\times\cJ\times\Gamma
\]
is then the unique least terminal-atomic coordinate subcarrier among
subsets of the literal carrier \(\cU\).

If instead the terminal criterion
\eqref{eq:completion-terminal-rectangularity-criterion} holds but the
pre-terminal criterion
\eqref{eq:completion-pre-rectangularity-criterion} fails, a pair relation
can represent the terminal atomic field exactly but does not preserve the
parent raw form.  A proof chain using that parent must retain the full
pre-terminal incidence or insert a separate proved parent comparison.
\end{theorem}

\begin{proof}
The first assertion follows from
\cref{thm:completion-least-terminal-faithful-relation}; unique coordinate
minimality follows from
\cref{prop:completion-minimal-coordinate-supports}.  The second assertion
is the simultaneous application of
\cref{thm:completion-least-terminal-faithful-relation,%
cor:completion-least-preterminal-faithful-relation}, together with the
positive discrepancy \eqref{eq:incidence-P-positive-difference}.
\end{proof}

The term \emph{hypergraph} records the four-label incidence; it introduces
no new averaging device.  Keeping it merely refuses an identification
which the pair projection cannot reconstruct.

\subsection{The noncircular fixed-shift audit}

Let \(\cU_X^{\mathrm{lit}}\) be obtained by copying every pre-terminal
coordinate of the upstream coefficientwise field, including its row,
cofactor, orbit, and opposite-row support tests, and put
\begin{equation}
 I_X^{\mathrm{lit}}:=\pr_{m,r}\cU_X^{\mathrm{lit}}.
 \label{eq:completion-literal-pair-projection}
\end{equation}
This projection alone does not prove that a later rectangular space is
faithful.

\begin{proposition}[Noncircular provenance criterion]
\label{prop:completion-noncircular-provenance-criterion}
Identification of a later retained relation \(\cI_X\) with the literal
pair projection is valid at the following respective levels only after
the corresponding statements have been proved:
\begin{align}
 \cI_X=I_X^{\mathrm{lit}}
 \quad\text{and}\quad
 \operatorname{PRect}(I_X^{\mathrm{lit}})
   &=(\cU_X^{\mathrm{lit}})^{\mathrm{pre},+},
 \label{eq:completion-required-pre-provenance}\\
 \cI_X=I_{X,h}^{\min}
 \quad\text{and}\quad
 \TRect_h(I_{X,h}^{\min})
   &=(\cU_X^{\mathrm{lit}})_h^+,
 \label{eq:completion-required-terminal-provenance}
\end{align}
where
\[
 I_{X,h}^{\min}:=\pr_{m,r}(\cU_X^{\mathrm{lit}})_h^+.
\]
The first line proves pre-terminal and hence terminal fidelity.  The second
proves terminal-atomic fidelity only.  Finiteness of \(\cI_X\), equality
of cardinalities, or agreement of dyadic cofactor ranges does not imply
either line.
\end{proposition}

\begin{proof}
Sufficiency and the exact logical strength of the two lines are
\cref{cor:completion-least-preterminal-faithful-relation,%
thm:completion-least-terminal-faithful-relation}.  Conversely, a missing
actual pair deletes a nonzero coordinate, while failed rectangularity adds
an active spurious cross.  Cardinality and range data record neither event,
as \cref{ex:completion-same-pair-different-incidence} demonstrates.
\end{proof}

\begin{remark}[No inherited identification is assumed]
\label{rem:completion-no-inherited-identification}
Neither \eqref{eq:completion-required-pre-provenance} nor
\eqref{eq:completion-required-terminal-provenance} is inferred from the
phrase \emph{supported pairs}.  In the physical application one must
return to the TPC--47 coefficient field, write
\(\cU_X^{\mathrm{lit}}\) by a literal set-builder formula, and verify the
chosen equality atom by atom.  If that verification fails, the correct
repair is the full-incidence carrier, not a redefinition of \(\cI_X\)
chosen to make exposure favorable.
\end{remark}

On the very-high face \(p>\sqrt{N_+}\), largest-prime factorization makes
the terminal cofactor unique for fixed \((m,j,\gamma)\).  This helps check
\eqref{eq:completion-required-terminal-provenance}, but does not prove the
required saturation: one must still show that every base label crossed
with a retained pair is present or terminal-null.
